\chapter{Gastro-Intestinal Illnesses}
\index{Gastro-Intestinal Illnesses}

\section*{Definition and Overview}
Gastro-intestinal (GI) illnesses are conditions affecting the digestive tract, which runs from the mouth through the stomach and intestines to the anus. They range from short-lived infections such as gastroenteritis, to chronic conditions such as irritable bowel syndrome, inflammatory bowel disease, coeliac disease, and reflux, to serious diseases including bowel cancer. GI illnesses are extremely common: most people experience nausea, vomiting, diarrhoea, constipation, or abdominal pain at some point, and chronic gut conditions affect a large proportion of the population. This chapter provides an overview of the main types of GI illness, their symptoms, and how they are diagnosed and treated.

\section*{Epidemiology}
GI illnesses are among the most common reasons for medical visits. Gastroenteritis affects millions of people each year, irritable bowel syndrome affects around 10--15% of the population, and reflux affects around 1 in 5 adults. Coeliac disease affects about 1 in 70 people, inflammatory bowel disease around 1 in 250, and bowel cancer is one of the most commonly diagnosed cancers, with around 15,000 new cases a year. Together, GI conditions cause substantial illness, lost work, and health costs, and screening for bowel cancer saves lives.

\section*{Aetiology and Pathophysiology}
\begin{itemize}
    \item \textbf{Infections:} viruses (norovirus, rotavirus), bacteria (salmonella, campylobacter, H.\ pylori), and parasites cause acute GI illness
    \item \textbf{Reflux disease:} stomach acid washes up into the oesophagus, damaging it, when the valve at the lower end fails
    \item \textbf{Peptic ulcers:} H.\ pylori and NSAIDs damage the stomach and duodenal lining
    \item \textbf{Irritable bowel syndrome:} a disorder of the gut-brain axis with altered motility and sensitivity, without visible inflammation
    \item \textbf{Inflammatory bowel disease:} ulcerative colitis and Crohn's disease are chronic immune inflammation of the bowel
    \item \textbf{Coeliac disease:} an autoimmune reaction to gluten that damages the small bowel lining
    \item \textbf{Bowel cancer:} develops from polyps over years, driven by age, genetics, and lifestyle
    \item \textbf{Functional disorders:} constipation, diarrhoea, and dyspepsia without a structural cause
    \item \textbf{Other conditions:} diverticular disease, gallstones, pancreatitis, and liver disease
\end{itemize}

\section*{Clinical Features}
\begin{itemize}
    \item Common symptoms: abdominal pain, bloating, nausea, vomiting, diarrhoea, constipation, reflux, and changes in bowel habit
    \item Acute illness: sudden vomiting and diarrhoea with gastroenteritis, usually settling within days
    \item Chronic pain and bloating with irritable bowel syndrome
    \item Bloody diarrhoea, weight loss, and fever with inflammatory bowel disease
    \item Weight loss, anaemia, and fatigue with coeliac disease and bowel cancer
    \item Heartburn and regurgitation with reflux
    \item Warning features: blood in the stool, black stools, unintended weight loss, and a change in bowel habit lasting more than a few weeks
    \item The pattern and duration distinguish acute from chronic illness
\end{itemize}

\section*{Differential Diagnosis}
The same symptoms can arise from many conditions.
\begin{itemize}
    \item Abdominal pain: infection, ulcers, gallstones, pancreatitis, appendicitis, and gynaecological causes
    \item Diarrhoea: infection, irritable bowel syndrome, inflammatory bowel disease, coeliac disease, and medicines
    \item Constipation: diet, medicines, endocrine disease, and obstruction
    \item Bleeding: haemorrhoids, polyps, inflammatory bowel disease, and cancer
    \item The history, examination, and targeted tests distinguish these
\end{itemize}

\section*{When to Seek Medical Care}
Most acute GI symptoms settle without treatment, but a doctor should be consulted for persistent symptoms, blood in the stool, weight loss, or a change in bowel habit, and for severe pain or vomiting.

\textbf{Contact your local emergency services for:}
\begin{itemize}
    \item Severe abdominal pain with vomiting and an inability to pass wind or stool
    \item Vomiting blood or passing black, tarry stools
    \item Significant bleeding from the back passage
    \item Signs of severe dehydration: no urine, confusion, or collapse
    \item Fever with severe abdominal pain
    \item Any emergency symptoms
\end{itemize}

\section*{Diagnosis and Investigations}
\begin{itemize}
    \item \textbf{History and examination:} the pattern of symptoms, warning features, and risk factors
    \item \textbf{Stool tests:} for infection, inflammation (calprotectin), and blood
    \item \textbf{Blood tests:} full blood count, inflammatory markers, coeliac serology, and liver and pancreatic function
    \item \textbf{Endoscopy:} gastroscopy and colonoscopy to examine the digestive tract directly and biopsy
    \item \textbf{Imaging:} ultrasound, CT, and MRI for structural and organ disease
    \item \textbf{Breath tests:} for H.\ pylori and lactose intolerance
    \item \textbf{Bowel cancer screening:} the national program's home stool test every 2 years from age 45--50
\end{itemize}

\section*{Management}
\begin{itemize}
    \item \textbf{Infections:} fluids and rest; antibiotics for specific bacterial infections
    \item \textbf{Reflux:} lifestyle measures, antacids, and acid-suppressing medicines
    \item \textbf{Irritable bowel syndrome:} dietary management (including the low-FODMAP diet under guidance), stress management, and medicines
    \item \textbf{Inflammatory bowel disease:} anti-inflammatory and immune-modulating medicines, with surgery for complications
    \item \textbf{Coeliac disease:} a strict lifelong gluten-free diet
    \item \textbf{Bowel cancer:} surgery, and chemotherapy or radiotherapy where needed
    \item \textbf{Constipation and diarrhoea:} fibre, fluids, activity, and treatment of the cause
    \item \textbf{Lifestyle:} diet, exercise, alcohol reduction, and smoking cessation
    \item \textbf{Follow-up and surveillance:} regular review for chronic conditions and after treatment
\end{itemize}

\section*{Complications}
\begin{itemize}
    \item Dehydration and electrolyte disturbance with severe diarrhoea and vomiting
    \item Ulcers bleeding and perforating
    \item Strictures, fistulas, and surgery in inflammatory bowel disease
    \item Malnutrition and anaemia in coeliac disease and bowel disease
    \item Bowel obstruction and spread of bowel cancer
    \item Reduced quality of life with chronic conditions
\end{itemize}

\section*{Prognosis and Outcomes}
Most GI illnesses resolve or are well managed. Gastroenteritis settles within days, reflux responds to treatment, irritable bowel syndrome improves with a tailored plan, and coeliac disease is controlled with a gluten-free diet. Inflammatory bowel disease is managed long term with good control in most people. Bowel cancer, when detected early through screening, has a high cure rate. The key to good outcomes is prompt assessment of warning features, accurate diagnosis, and ongoing management of chronic conditions.

\section*{Prevention and Self-Care}
\begin{itemize}
    \item Do the National Bowel Cancer Screening test when it arrives
    \item Follow food safety rules to prevent gastroenteritis
    \item Eat a high-fibre diet, stay active, and maintain a healthy weight
    \item Limit alcohol, stop smoking, and avoid excessive NSAID use
    \item Manage stress, which affects gut symptoms
    \item See a doctor early for blood in the stool, weight loss, or a change in bowel habit
    \item Follow treatment and surveillance plans for chronic conditions
\end{itemize}

\section*{Sources and Further Reading}
The following reputable sources provide further information for healthcare students and patients seeking reliable, current advice about gastro-intestinal illnesses.
\begin{itemize}
    \item The Gut Foundation. \textit{Digestive health information.} https://www.gutfoundation.org.au/
    \item Healthdirect Australia. \textit{Digestive health.} https://www.healthdirect.gov.au/digestive-system
    \item Gastroenterology Society of Australia (GESA). \textit{Patient information.} https://www.gesa.org.au/
    \item Bowel Cancer Australia. \textit{Bowel health and screening.} https://www.bowelcanceraustralia.org/
    \item NHS. \textit{Digestive health.} https://www.nhs.uk/
    \item World Health Organization. \textit{Digestive diseases.} https://www.who.int/
\end{itemize}
